\clearpage\thispagestyle{empty}
\pdfbookmark[0]{Quellen zum Ceres-Bericht vom September und Oktober 1803}{ceres18030910-sources}
\begin{center}\textbf{Quellen zum Ceres-Bericht vom September und Oktober 1803}\end{center}
\noindent Diese KI-Übertragung in heutiges Deutsch folgt dem Auszug aus dem Bericht von Franz Xaver von Zach in \href{https://gdz.sub.uni-goettingen.de/id/PPN236007467}{Gauß’ \textit{Werken}, Band~VI} (Göttingen, 1874), S.~245--246. Der Abschnitt über die verbesserten Elemente und die beiden Tabellen entsprechen \href{https://api.digitale-sammlungen.de/iiif/image/v2/bsb10538600_00298/full/1450,/0/default.jpg}{S.~290} des Beitrags XXXII, \href{https://api.digitale-sammlungen.de/iiif/image/v2/bsb10538600_00296/full/1450,/0/default.jpg}{S.~288--291}, im Septemberheft der \textit{Monatlichen Correspondenz}, Band~VIII (1803). Die drei Beobachtungen von Giuseppe Piazzi stehen im \href{https://api.digitale-sammlungen.de/iiif/image/v2/bsb10538600_00100/full/1450,/0/default.jpg}{Juliheft, S.~94}. Die Tabelle der \textit{Werke} gibt die mittleren Zeiten in Palermo an und bezeichnet die Unterschiede mit R—B.
\par\smallskip\noindent Der Absatz in eckigen Klammern in den \textit{Werken} beschreibt die von Carl Friedrich Gauß berechnete und im Beitrag XLIII, \href{https://api.digitale-sammlungen.de/iiif/image/v2/bsb10538600_00379/full/1450,/0/default.jpg}{S.~369--371} des Oktoberhefts, gedruckte Ephemeride. Sie steht auf \href{https://api.digitale-sammlungen.de/iiif/image/v2/bsb10538600_00380/full/1450,/0/default.jpg}{S.~370} und \href{https://api.digitale-sammlungen.de/iiif/image/v2/bsb10538600_00381/full/1450,/0/default.jpg}{S.~371}. Clifford Cunningham gibt den Oktoberbeitrag in \href{https://doi.org/10.1007/978-3-319-28815-4_10}{„The Scientific Papers on Ceres“} (2016), S.~364--365, auf Englisch wieder und zeigt die Ephemeride in Abb.~10.68 und~10.69.
\par\smallskip\noindent Übertragung und Satz wurden mit KI erstellt. Die Digitalisate der Zeitschrift stammen von der Bayerischen Staatsbibliothek.
\clearpage
